\chapter{The Computational Companion}\label{app:companion}

The directory \texttt{companion/} holds a finite-state implementation in Python using only the standard library. Its output is recorded in \texttt{companion/EXPECTED\_OUTPUT.txt}. The checks below are \emph{reported checks}: they have been run by the author of the skeleton, and the package and recorded output should be inspected before any of them is cited as verified. Run \texttt{python3 -m vcc.run all} from \texttt{companion/}, or \texttt{python3 -m unittest discover tests}.

Three kinds of check are distinguished. An \emph{exhaustive} check enumerates every case of a small model and so establishes the claim for that model only. A \emph{random} check samples systems from a recorded seed. It can find counterexamples and give implementation evidence, but it cannot prove a universal theorem. A \emph{counterexample} check exhibits a specific system in which a hypothesis has been removed.

\begin{center}
\small
\renewcommand{\arraystretch}{1.2}
\begin{tabular}{@{}llllll@{}}
\toprule
Result & Claim checked & Kind & Seed & Hypothesis removed\\
\midrule
\Pref{P03} & iteration = Knaster--Tarski union & random, 300 systems & 1 & --\\
\Pref{P03} & no termination on an infinite chain & counterexample & -- & finiteness\\
\Pref{P05} & local acceptability without viability & exhaustive, 3 states & -- & --\\
\Pref{P09} & each demand fits, the sum does not & random, 500 draws & 2 & --\\
\Pref{P10} & $K_{12}=K_1\times K_2$ & random, 200 pairs & 3 & --\\
\Pref{P10} & $K_{12}\subsetneq K_1\times K_2$ & counterexample & -- & independence\\
\Pref{P23} & monotone in $G$, actions, disturbances & random, 300 systems & 4 & --\\
\Pref{P22} & potential vs.\ negative cycle & exhaustive, 3 nodes & -- & friction varied\\
\Pref{P28} & lost-exit interval in closed form & random, 400 parameter sets & 5 & --\\
\Pref{P29} & availability, timing, sufficiency & random, 600 cases & 6 & --\\
\bottomrule
\end{tabular}
\end{center}

Random checks confirm closed forms that are proved in the text; they are not part of any proof. The housing-case numbers in Theorems~H1 and H2 are reproduced exactly by the last two rows.
